\documentclass{beamer}
\usepackage{amsmath}
\setbeameroption{show notes on second screen=right}
\setbeamertemplate{note page}[plain]

\begin{document}

\begin{frame}
\frametitle{}


$\frac35 - \alert<3>{\frac25 \times \frac56}$. \note[item]{\alert<1>{Find three fifths minus two fifths times five sixths.}}

\pause\note[item]{\alert<2>{It is tempting to do subtraction first because it is written first,}}
\pause\note[item]{\alert<3>{but we have to focus on multiplying first because of the order of operations. So, first let's simplify two fifths times five sixths.}}

\vskip4pt

\pause $= \frac{3}{5} - \frac{2 \times 5}{5 \times 6}$
\note[item]{\alert<4>{For multiplying fractions, we multiply straight across, so we have 2 times 5 over 5 times 6. Even though the first fraction 3 over 5 has not changed, we wrote it right after the equal sign right before the minus sign to account for it.}}

\vskip4pt

\pause $= \frac{3}{5} - \frac{10}{30}$
\note[item]{\alert<5>{After copying down the first fraction 3 over 5, we can either rewrite the second fraction as 10 over 30.}}

\vskip4pt

\pause $= \frac{3}{5} - \frac{1}{3}$
\note[item]{\alert<6>{After copying down the first fraction 3 over 5, we can also simplify the second fraction. Since 10 and 30 share a common factor of 10, we can divide the numerator and denominator by 10 to get 1 over 3.}}

\vskip4pt

\pause $= \frac{9}{15} - \frac{5}{15}$
\note[item]{\alert<7>{We rewrite each fraction with a denominator of 15. By multiplying the top and bottom of the first fraction by 3, we replace the first fraction with 9 over 15.  By multiplying the top and bottom of the second fraction by 5, we replace the second fraction with 5 over 15.}}

\vskip4pt

\pause $= \frac{9-5}{15}$
\note[item]{\alert<8>{We subtract the numerators while keeping the denominator the same. So we have 9 minus 5 over 15.}}

\vskip4pt

\pause $= \frac{4}{15}$
\note[item]{\alert<9>{By simplifying the numerator, our final answer is 4 over 15. This fraction is already in its simplest form since 4 and 15 do not share any common factors.}}


\end{frame}



%% Blank frames at the end to prevent weird shifts.
\begin{frame}
\end{frame}
\begin{frame}
\end{frame}

\end{document}
